% !TEX root = ../frode-bibelen.tex
% ----------------------------------------------------------------------------
%  Frode-bibelen -- 12 Eksamensarkivet H19--S26
%  Inkluderes fra frode-bibelen.tex (kompiler alltid hovedfilen).
% ----------------------------------------------------------------------------

% ================================================== 12 EKSAMENSARKIV =======
\kapittel{Eksamensarkivet H19--S26}

Alle eksamenene: 4 timer, oppgave 1 = 30\,\%, oppgave 2 = 70\,\%. Oppgave 2 er \emph{alltid} et menystyrt program med
global array av pekere (eller structer) + teller, og nesten alltid \cc{lesFraFil} som siste deloppgave.
Lenkene går til \href{https://frh.folk.ntnu.no/grprog/eksamen/}{\cc{grprog/eksamen/}} (alt finnes også i \cc{EKSAMEN.zip}).

\begin{center}
\small
\begin{longtable}{@{}L{0.9cm}L{4.0cm}L{8.6cm}L{1.8cm}@{}}
\toprule
\textbf{Eks.} & \textbf{Oppgave 2 handler om} & \textbf{Datastruktur / spesielt} & \textbf{Filer} \\
\midrule
\endhead
\textbf{S26} & Kunder som bestiller vindusglass fra ulike produsenter & Pekere \cc{gKundene[]} + array \emph{av} structer \cc{gGlassene[]}; \cc{int glass[MAXGLASS]} teller per glasstype; kvm og faktura (\cc{float}); to datafiler. & \href{https://frh.folk.ntnu.no/grprog/eksamen/ex_s26.pdf}{oppgave} \newline \href{https://frh.folk.ntnu.no/grprog/eksamen/ex_s26_2.c}{løsning 2} \newline \href{https://frh.folk.ntnu.no/grprog/eksamen/ex_s26_2.tpl}{skjelett} \\
\textbf{H25} & Personer i en sosial gruppe og dager de kan møtes & \cc{bool dag[31]} i structen; fjern deltager (flytt siste inn); dager der alle/minst halvparten kan; 0/1 fra fil til \cc{bool}. & \href{https://frh.folk.ntnu.no/grprog/eksamen/ex_h25.pdf}{oppgave} \newline \href{https://frh.folk.ntnu.no/grprog/eksamen/ex_h25_2.c}{løsning 2} \newline \href{https://frh.folk.ntnu.no/grprog/eksamen/ex_h25_2.tpl}{skjelett} \\
\textbf{S25} & Besøk på en legevakt per ukedag og time & Fast \cc{gDogn[7]} med 24 timer hver; når det er flest besøk; skriv til og les fra fil. & \href{https://frh.folk.ntnu.no/grprog/eksamen/ex_s25.pdf}{oppgave} \newline \href{https://frh.folk.ntnu.no/grprog/eksamen/ex_s25_2.c}{løsning 2} \newline \href{https://frh.folk.ntnu.no/grprog/eksamen/ex_s25_2.tpl}{skjelett} \\
\textbf{H24} & Havner og båtrutene som går innom dem & To structer (\cc{Havn}, \cc{Rute}); en rute har opptil 20 anløp; mest travle havn via telle-array; to datafiler. & \href{https://frh.folk.ntnu.no/grprog/eksamen/ex_h24.pdf}{oppgave} \newline \href{https://frh.folk.ntnu.no/grprog/eksamen/ex_h24_2.c}{løsning 2} \newline \href{https://frh.folk.ntnu.no/grprog/eksamen/ex_h24_2.tpl}{skjelett} \\
\textbf{S24} & Én klubbs kamper og egne målscorere & Referanseoverføring \cc{kampHent(kamp, \&for, \&mot)}; poengsum og «tabellinje»; toppscorer(e). & \href{https://frh.folk.ntnu.no/grprog/eksamen/ex_s24.pdf}{oppgave} \newline \href{https://frh.folk.ntnu.no/grprog/eksamen/ex_s24_2.c}{løsning 2} \newline \href{https://frh.folk.ntnu.no/grprog/eksamen/ex_s24_2.tpl}{skjelett} \\
\textbf{H23} & Fuglearter og observasjonene av dem & \cc{antObs} + \cc{observasjoner[365]} (\cc{MMDD}) i structen; \cc{finnFugl} på navn; les fil til \cc{feof}; dag med flest observasjoner. & \href{https://frh.folk.ntnu.no/grprog/eksamen/ex_h23.pdf}{oppgave} \newline \href{https://frh.folk.ntnu.no/grprog/eksamen/ex_h23_2.c}{løsning 2} \newline \href{https://frh.folk.ntnu.no/grprog/eksamen/ex_h23_2.tpl}{skjelett} \\
\textbf{S23} & Brøytere/strøere og kundene deres & To structer; registrere brøyting/strøing; lage fakturaer; hvem som tjener mest. & \href{https://frh.folk.ntnu.no/grprog/eksamen/ex_s23.pdf}{oppgave} \newline \href{https://frh.folk.ntnu.no/grprog/eksamen/ex_s23_2.c}{løsning 2} \newline \href{https://frh.folk.ntnu.no/grprog/eksamen/ex_s23_2.tpl}{skjelett} \\
\textbf{H22} & Hendelser på en biltur i Alpene & \cc{enum Type} med 6 verdier; dato \cc{MMDD} og klokkeslett \cc{TTMM}; dato med flest hendelser. & \href{https://frh.folk.ntnu.no/grprog/eksamen/ex_h22.pdf}{oppgave} \newline \href{https://frh.folk.ntnu.no/grprog/eksamen/ex_h22_2.c}{løsning 2} \newline \href{https://frh.folk.ntnu.no/grprog/eksamen/ex_h22_2.tpl}{skjelett} \\
\textbf{S22} & Ting som må handles ved oppussing & \cc{enum Enhet}; søk på tekst; kostnad (\cc{float}) per ting. & \href{https://frh.folk.ntnu.no/grprog/eksamen/ex_s22.pdf}{oppgave} \newline \href{https://frh.folk.ntnu.no/grprog/eksamen/ex_s22_2.c}{løsning 2} \newline \href{https://frh.folk.ntnu.no/grprog/eksamen/ex_s22_2.tpl}{skjelett} \\
\textbf{H21} & Kunders lån i ulike banker & To structer (\cc{Kunde}, \cc{Bank}); \cc{bankNr[2]} som indekser inn i \cc{gBanker}; renter (\cc{float}); to datafiler. & \href{https://frh.folk.ntnu.no/grprog/eksamen/ex_h21.pdf}{oppgave} \newline \href{https://frh.folk.ntnu.no/grprog/eksamen/ex_h21_2.c}{løsning 2} \newline \href{https://frh.folk.ntnu.no/grprog/eksamen/ex_h21_2.tpl}{skjelett} \\
\textbf{S21} & Rakfisk- og pultostprodukter og produsentene & \cc{enum Type}; produsentnavn i global tekst-array; \cc{strstr}-søk; skriv til fil (\cc{.DT2}). & \href{https://frh.folk.ntnu.no/grprog/eksamen/ex_s21.pdf}{oppgave} \newline \href{https://frh.folk.ntnu.no/grprog/eksamen/ex_s21_2.c}{løsning 2} \newline \href{https://frh.folk.ntnu.no/grprog/eksamen/ex_s21_2.tpl}{skjelett} \\
\textbf{H20} & Gjenstander som skal ryddes ut av et hus & \cc{enum Kategori} med 6 verdier; skriv alle i én kategori; endre kategori; skriv til og les fra fil. & \href{https://frh.folk.ntnu.no/grprog/eksamen/ex_h20.pdf}{oppgave} \newline \href{https://frh.folk.ntnu.no/grprog/eksamen/ex_h20_2.c}{løsning 2} \newline \href{https://frh.folk.ntnu.no/grprog/eksamen/ex_h20_2.tpl}{skjelett} \\
\textbf{S20} & Episoder sett i sesongene av en lang serie & Episoder per sesong i structen; søk på sesongnavn (egen \cc{bool}-funksjon); prosentvis sett. & \href{https://frh.folk.ntnu.no/grprog/eksamen/ex_s20.pdf}{oppgave} \newline \href{https://frh.folk.ntnu.no/grprog/eksamen/ex_s20_2.c}{løsning 2} \newline \href{https://frh.folk.ntnu.no/grprog/eksamen/ex_s20_2.tpl}{skjelett} \\
\textbf{H19} & Ansatte, arbeidsdager og oppmøte & Dager i måneden per ansatt; registrere oppmøte; statistikk; to datafiler (ansatte og jobbing). & \href{https://frh.folk.ntnu.no/grprog/eksamen/ex_h19.pdf}{oppgave} \newline \href{https://frh.folk.ntnu.no/grprog/eksamen/ex_h19_2.c}{løsning 2} \newline \href{https://frh.folk.ntnu.no/grprog/eksamen/ex_h19_2.tpl}{skjelett} \\
\bottomrule
\end{longtable}
\end{center}

\begin{tips}[{Slik øver du med arkivet}]
\begin{enumerate}
  \item Last ned \cc{ex\_xxx\_2.tpl} (skjelettet) og datafilene (\cc{.dta} -- gi dem filnavnet oppgaven bruker, f.eks.\ \cc{EX\_H25\_SAMLING.DTA}).
  \item Løs oppgaven på 2,5 timer med \emph{bare} papir-hjelpemidler (denne bibelen + \cc{LesData.h}).
  \item Kompiler og test. Sammenlign med \cc{ex\_xxx\_2.c} og rettemalen -- skriv det du bommet på i margen her.
\end{enumerate}
\end{tips}
